\section{Experimental design}\label{sec:design}

Table~\ref{tab:params} lists the parameters of the base design. Two pre-registered ablations change one learner parameter each ($\gamma = 0$ and memory~0, Section~\ref{sec:indicators}); the Phase~1 experiments each change one element of the design (the horizon following where the exploration schedule requires) and are specified in Section~\ref{sec:phase1}. The Market and Learner parameters, the 1M horizon and the convergence window come from the project specification; the seeds per cell, the frozen-play length, its burn-in and the CI rule are project choices; all were fixed before any result was seen and none was tuned. The one base parameter added after a result was seen is the 5M horizon (``Two horizons'' below). Apart from $\beta$ and the horizon no hyperparameter was varied in search of a result; the constant-$\varepsilon$ levels of Section~\ref{sec:occupation} and the two restart levels $\varepsilon_0$ of Section~\ref{sec:reinject} are the independent variables of those experiments, fixed with their predictions.

\begin{table}[htbp]
\centering
\caption{Parameters of the base design. Notes: the Market and Learner groups, the 1M horizon and the convergence window come from the specification; the seeds per cell, frozen-play length, burn-in and CI rule are design choices fixed before any result was seen; the 5M horizon was added after the 1M results (see ``Two horizons''); the punishment windows and threshold and the collapse criterion are discussed under ``What was fixed in advance''; ``rounds'' means auction rounds; SE is the standard error across seeds.}
\label{tab:params}
\begin{tabular}{llr}
\toprule
Group & Parameter & Value \\
\midrule
Market & Firms & 3 \\
 & Capacity per firm & 60 MW \\
 & Marginal cost & \eur{10}/MWh \\
 & Demand (inelastic) & 100 MW \\
 & Price cap & \eur{100}/MWh \\
 & Bid grid & 15 rungs, \eur{10} to \eur{100} \\
 & Grid rung & \eur{6.43} \\
Learner & Learning rate $\alpha$ & 0.15 \\
 & Discount $\gamma$ & 0.95 (ablation: 0) \\
 & Exploration decay $\beta$ & $10^{-5}$ (sweep: $10^{-6}$ to $10^{-4}$) \\
 & Memory $m$ & 1 (ablation and control: 0) \\
 & States, $m=1$ & 3,375 \\
 & Initial table $Q_0$ & $\mathbb{E}[\text{profit}(a,\ \text{rivals uniform})]/(1-\gamma)$ \\
Run & Horizon $T$ & 1M and 5M rounds \\
 & Convergence window & 100,000 rounds \\
 & Frozen play & 1,000 rounds \\
 & Frozen burn-in dropped & 100 rounds \\
 & Seeds per cell & 100 \\
 & CI rule & mean $\pm 1.96 \times$ SE \\
Punishment test & Settling rounds & 200 \\
 & Pre / post rounds & 6 / 24 \\
 & Judgement windows & 2 $\times$ 12 rounds \\
 & Reaction threshold & half a rung (\eur{3.21}) \\
Ablation & Collapse criterion & $\Dl < 0.25$ and $< \tfrac{1}{2}$ full learner \\
\bottomrule
\end{tabular}
\end{table}

\paragraph{What was fixed in advance, and what was not.}
``Pre-registered'' here means written down in the project's own documents before the corresponding runs; there is no public registry entry, and the record is a git history rather than a timestamped filing. Three levels should be distinguished, because they do not carry equal weight.

The specification, written before any code, fixes the market, the learners, every hyperparameter in the Market and Learner groups of Table~\ref{tab:params}, the 1M horizon, the four measurements, and acceptance rules stated in words: a meaningful fraction of seeds above $\Dl = 0.5$ for M2; prices that drop below the pre-deviation level and then recover for M3; $\Dl$ collapsing toward zero under both ablations for M4; at least one strictly profitable unplayed deviation for M5. None of these carries a number.

The numerical thresholds were therefore fixed one level down, in the implementation, and are present in the first committed version of the code, before any run reported here. They are the collapse rule for M4 ($\Dl < 0.25$ and below half the full learner's $\Dl$) and the whole operationalisation of M3: judging on the rivals' bids against a no-deviation counterfactual, in two 12-round windows, with a half-rung threshold. Two of these choices were informed by earlier runs and we flag them rather than claim more than the record supports. The 12-round window length is justified by the converged cycle lengths, which were known only after the first runs: it lets a cycle of period 1, 2, 3, 4 or 6 complete a whole number of times in each window, while periods 5 and 7, which occur in 9 of the 100 long-run seeds, are not covered and can be misread (Section~\ref{sec:discussion}); at 1M rounds, where periods run up to 56 and only 23 of 100 seeds have period four or less, the 12-seed verdicts of Table~\ref{tab:m3} are more exposed to this than the long-run ones. And the rival-bid criterion replaced an earlier price-based one after the price-based version was found to read a converged cycle's own price dip as punishment. Every M3 verdict reported here was computed with the final criterion, and no verdict computed under the earlier one is reported.

The Phase~1 experiments are not in the specification at all. They were designed after the results of Section~\ref{sec:results} were seen, and sequentially, each in response to the last. What was fixed in advance there is narrower: each experiment's prediction was written into the script before that experiment ran, and all five are reported with their outcomes, two of which failed outright, in Section~\ref{sec:phase1}.

\paragraph{Acceptance rules.}
No numerical acceptance threshold for M2 is on record: the specification asks for a meaningful fraction of seeds that converge to $\Dl > 0.5$. We report the share of seeds above $\Dl = 0.5$ and call M2 met only where every seed exceeds it, which is the case at 5M rounds (100\%, all converged) and not at 1M (12\%, none converged); readers who accept a lower threshold may read the 1M result differently. M3 is judged with firm~1 as the forced deviator, the code's default; the verdicts for firms~2 and~3 are reported in Appendix~\ref{app:extra}, and no rule for combining the three was fixed in advance. M4 is met only if both ablated learners collapse, $\Dl < 0.25$ and below half the full learner's $\Dl$. M5 is a necessary condition: the collusion reading requires an unclaimed gap, and the memory-0 control tests whether it is sufficient. The outcome is read as supporting tacit collusion only if M2, M3 (punish--forgive) and M4 are all met and M5 fires for the memory-1 learner but not for the control.

\paragraph{Seeds and confidence intervals.}
Every configuration is run on 100 independent seeds; where a table reports a subset (the seeds with $\Dl > 0.5$, labelled ``collusive'' in Section~\ref{sec:results}) it says so. A seed fixes the random draws of a run; the initial tables are deterministic. Means are reported with a 95\% confidence interval of mean $\pm 1.96$ times the standard error across seeds; medians, ranges and per-seed counts are reported with the price levels in Section~\ref{sec:results}. A 50-seed run of the whole study preceded the 100-seed run, on the first 50 of the same seeds; every $\Dl$ moved by at most 0.02 between the two, and no verdict changed.

\paragraph{Two horizons.}
At the specified $T = 1$M rounds the learners are almost greedy ($\varepsilon_t = 4.5 \times 10^{-5}$, Section~\ref{sec:learners}), yet $\Dl = 0.375$, no seed meets the convergence criterion and M2 is not met (Section~\ref{sec:results}). The 5M-round horizon is not in the specification: it was added after that result was seen and is not pre-registered, and every base configuration converges in every seed there. It is one of two cells chosen after seeing the quantity it reports; the other is the 30M-round $\varepsilon = 10^{-4}$ point of Table~\ref{tab:occupation}, added once the memory-1 occupation curve was seen still rising at $\varepsilon = 10^{-3}$ (Section~\ref{sec:occupation}). We report both horizons: M2 fails at 1M and is met at 5M, M4 and M5 give the same verdict at both, and M3 is judged on 12 seeds at 1M and on all 100 at 5M; the headline $\Dl = 0.621$ comes from the added one. Every run goes to its full horizon.

\paragraph{What is measured, and on what.}
$\Dl$ and all behavioural indicators are computed on frozen play as defined in Section~\ref{sec:indicators}: 1,000 greedy rounds from the final state, the first 100 dropped. Outside the constant-$\varepsilon$ experiment of Section~\ref{sec:occupation}, no reported $\Dl$ or verdict is computed on the training trace, because it mixes exploration draws with the learned policy in a share that changes with $t$; the training-trace prices in Figures~\ref{fig:trajectory} and~\ref{fig:trajectory5m} and the lowest 1,000-round mean of Table~\ref{tab:reinject} are descriptive only. Frozen play isolates the greedy policy from the final state; it visits only the cycle reachable from that state, and M5 and the Bellman residual are computed on that cycle alone. The punishment test adapts the one-round forced deviation of \citet{calvano2020}, with the verdict read on the rivals' bids against a no-deviation counterfactual because, from any profile at which the two lowest bids are tied, a lone undercut does not move the price (Section~\ref{sec:market}), so the price is an unreliable signal of a reaction.

\paragraph{The memory-0 control.}
The memory-0 learner sees one state, so it cannot condition on what rivals did last round and cannot carry a threat; on the working definition of Section~\ref{sec:related}, a firm holding back because it expects retaliation \citep{tirole1988,ivaldi2003}, whatever price it reaches is by construction not tacit collusion. In a Q-learner such a threat, if it exists, lives in the learned continuation values, and the control calibrates the indicators against that. An indicator that reads the same on both learners does not by itself distinguish a learned continuation from a stale estimate; and an indicator that separates them because a one-state table has nothing to react to shows that the memory-1 policy conditions on rivals' bids, which its state permits but does not guarantee, and not that the conditioning is a threat (Section~\ref{sec:discussion}).

\paragraph{Compute and reproducibility.}
The auction and learning loop are compiled with numba and run at about 0.4~s per 1M rounds per seed; the six base configurations at 100 seeds take about 9 minutes and the whole study about two hours sequentially on a laptop with no GPU. The code is the Python package \texttt{collusim} v0.1.0 with 31 tests (Appendix~\ref{app:repro}); every simulation number in this paper is taken from the stored \texttt{results/*.json} files, which the scripts described there produce, and the one-shot payoffs, Nash equilibria and deviation gains of Section~\ref{sec:market} are computed from the payoff table with the same package.
